\documentclass[aspectratio=169,10pt]{beamer}
\setbeamertemplate{navigation symbols}{}
\setbeamertemplate{footline}{\hfill\scriptsize\insertframenumber/\inserttotalframenumber\hspace*{3mm}\vspace*{1.5mm}}
\setbeamerfont{frametitle}{size=\large,series=\bfseries}
\setbeamercolor{frametitle}{fg=blue!40!black}
\setbeamertemplate{itemize items}{\textbullet}
\usepackage{booktabs,amsmath,amssymb,graphicx,xcolor,tabularx}
\newcommand{\yes}{\textcolor{green!45!black}{\textbf{yes}}}
\newcommand{\no}{\textcolor{red!70!black}{\textbf{no}}}
\newcommand{\partly}{\textcolor{orange!80!black}{\textbf{partly}}}
\newcommand{\verdict}[1]{\colorbox{gray!13}{\parbox{\dimexpr\linewidth-2\fboxsep}{\scriptsize\textbf{Verdict.} #1}}}
\newcommand{\plain}[2]{\fcolorbox{blue!35!black}{blue!4}{\parbox{\dimexpr\linewidth-2\fboxsep-2\fboxrule}{\footnotesize\textbf{#1}\enspace #2}}}
\newcommand{\why}[1]{\fcolorbox{blue!35!black}{blue!4}{\parbox{\dimexpr\linewidth-2\fboxsep-2\fboxrule}{\scriptsize\textbf{Why, in plain words.}\enspace #1}}}
\newcommand{\caseframe}[6]{%
\begin{frame}{#1}
\begin{columns}[T]
\begin{column}{0.53\textwidth}
\includegraphics[width=\linewidth,height=0.56\textheight,keepaspectratio]{../figs/#2}\\[0.5mm]
\resizebox{\linewidth}{!}{\input{#4}}
\end{column}
\begin{column}{0.45\textwidth}
\scriptsize
\textbf{Known optimum.} #3\par\smallskip
\verdict{#5}\par\smallskip
\why{#6}
\end{column}
\end{columns}
\end{frame}}
\newcommand{\ppm}{\,ppm}
% generated by tests/t18_numbers.py from out/t17_planar.json -- do not edit
\newcommand{\PlanarPairJ}{0.01}
\newcommand{\PlanarPairJwalk}{0.00}
\newcommand{\PlanarPairWall}{0.03}
\newcommand{\PlanarYard}{1.89}
\newcommand{\PlanarBWalk}{-9.15}
\newcommand{\PlanarBFine}{-10.3}
\newcommand{\PlanarBLim}{-11.9}
\newcommand{\PlanarBMinusFitWalk}{+1.70}
\newcommand{\PlanarBMinusFitLim}{-0.99}
\newcommand{\PlanarBTrials}{5}
\newcommand{\PlanarBRecords}{5}
\newcommand{\PlanarBTime}{100\,s}
\newcommand{\PlanarBStop}{converged}
\newcommand{\PlanarBPred}{0.003}
\newcommand{\PlanarCWalk}{-9.15}
\newcommand{\PlanarCFine}{-10.3}
\newcommand{\PlanarCLim}{-11.9}
\newcommand{\PlanarCMinusFitWalk}{+1.71}
\newcommand{\PlanarCMinusFitLim}{-0.99}
\newcommand{\PlanarCTrials}{13}
\newcommand{\PlanarCRecords}{13}
\newcommand{\PlanarCTime}{3\,min}
\newcommand{\PlanarCStop}{converged}
\newcommand{\PlanarCPred}{2e-05}
\newcommand{\PlanarAOneWalk}{-9.15}
\newcommand{\PlanarAOneFine}{-10.3}
\newcommand{\PlanarAOneLim}{-11.9}
\newcommand{\PlanarAOneMinusFitWalk}{+1.71}
\newcommand{\PlanarAOneMinusFitLim}{-0.98}
\newcommand{\PlanarAOneTrials}{11}
\newcommand{\PlanarAOneRecords}{11}
\newcommand{\PlanarAOneTime}{3\,min}
\newcommand{\PlanarAOneStop}{converged}
\newcommand{\PlanarAOnePred}{0.0002}
\newcommand{\PlanarBOneWalk}{-9.15}
\newcommand{\PlanarBOneFine}{-10.3}
\newcommand{\PlanarBOneLim}{-11.9}
\newcommand{\PlanarBOneMinusFitWalk}{+1.70}
\newcommand{\PlanarBOneMinusFitLim}{-0.99}
\newcommand{\PlanarBOneTrials}{5}
\newcommand{\PlanarBOneRecords}{5}
\newcommand{\PlanarBOneTime}{111\,s}
\newcommand{\PlanarBOneStop}{converged}
\newcommand{\PlanarBOnePred}{0.003}
\newcommand{\PlanarFitWalk}{-10.9}
\newcommand{\PlanarFitFine}{-8.96}
\newcommand{\PlanarFitLim}{-10.9}
\newcommand{\PlanarFitMinusFitWalk}{+0.00}
\newcommand{\PlanarFitMinusFitLim}{+0.00}
% generated by tests/t18_numbers.py from out/t17_methane.json -- do not edit
\newcommand{\MethanePairJ}{0.07}
\newcommand{\MethanePairJwalk}{0.00}
\newcommand{\MethanePairWall}{0.01}
\newcommand{\MethaneYard}{0.70}
\newcommand{\MethaneAWalk}{+1.01}
\newcommand{\MethaneAFine}{-2.11}
\newcommand{\MethaneALim}{-2.12}
\newcommand{\MethaneAMinusFitWalk}{+1.95}
\newcommand{\MethaneAMinusFitLim}{-1.87}
\newcommand{\MethaneATrials}{6}
\newcommand{\MethaneARecords}{6}
\newcommand{\MethaneATime}{52\,min}
\newcommand{\MethaneAStop}{converged}
\newcommand{\MethaneAPred}{0.001}
\newcommand{\MethaneCWalk}{+1.01}
\newcommand{\MethaneCFine}{-2.18}
\newcommand{\MethaneCLim}{-2.19}
\newcommand{\MethaneCMinusFitWalk}{+1.95}
\newcommand{\MethaneCMinusFitLim}{-1.95}
\newcommand{\MethaneCTrials}{8}
\newcommand{\MethaneCRecords}{8}
\newcommand{\MethaneCTime}{101\,min}
\newcommand{\MethaneCStop}{converged}
\newcommand{\MethaneCPred}{0.002}
\newcommand{\MethaneFitWalk}{-0.94}
\newcommand{\MethaneFitFine}{-0.24}
\newcommand{\MethaneFitLim}{-0.30}
\newcommand{\MethaneFitMinusFitWalk}{+0.00}
\newcommand{\MethaneFitMinusFitLim}{+0.00}
% generated by tests/t18_numbers.py from out/t17_rao61.json -- do not edit
\newcommand{\RaoSixtyOnePairJ}{0.05}
\newcommand{\RaoSixtyOnePairJwalk}{0.00}
\newcommand{\RaoSixtyOnePairWall}{0.03}
\newcommand{\RaoSixtyOneYard}{0.87}
\newcommand{\RaoSixtyOneBumpWalk}{-0.52}
\newcommand{\RaoSixtyOneBumpFine}{-1.30}
\newcommand{\RaoSixtyOneBumpLim}{-1.31}
\newcommand{\RaoSixtyOneBumpMinusFitWalk}{+0.77}
\newcommand{\RaoSixtyOneBumpMinusFitLim}{-0.97}
\newcommand{\RaoSixtyOneBumpTrials}{4}
\newcommand{\RaoSixtyOneBumpRecords}{4}
\newcommand{\RaoSixtyOneBumpTime}{71\,min}
\newcommand{\RaoSixtyOneBumpStop}{converged}
\newcommand{\RaoSixtyOneBumpPred}{0.0002}
\newcommand{\RaoSixtyOneHalfWalk}{-0.53}
\newcommand{\RaoSixtyOneHalfFine}{-1.35}
\newcommand{\RaoSixtyOneHalfLim}{-1.36}
\newcommand{\RaoSixtyOneHalfMinusFitWalk}{+0.77}
\newcommand{\RaoSixtyOneHalfMinusFitLim}{-1.02}
\newcommand{\RaoSixtyOneHalfTrials}{8}
\newcommand{\RaoSixtyOneHalfRecords}{8}
\newcommand{\RaoSixtyOneHalfTime}{100\,min}
\newcommand{\RaoSixtyOneHalfStop}{converged}
\newcommand{\RaoSixtyOneHalfPred}{0.002}
\newcommand{\RaoSixtyOneOldBumpWalk}{-8.29}
\newcommand{\RaoSixtyOneOldBumpFine}{-8.87}
\newcommand{\RaoSixtyOneOldBumpLim}{-8.64}
\newcommand{\RaoSixtyOneOldBumpMinusFitWalk}{-6.99}
\newcommand{\RaoSixtyOneOldBumpMinusFitLim}{-8.45}
\newcommand{\RaoSixtyOneOldHalfWalk}{-13.4}
\newcommand{\RaoSixtyOneOldHalfFine}{-14.3}
\newcommand{\RaoSixtyOneOldHalfLim}{-14.2}
\newcommand{\RaoSixtyOneOldHalfMinusFitWalk}{-12.1}
\newcommand{\RaoSixtyOneOldHalfMinusFitLim}{-13.9}
\newcommand{\RaoSixtyOneFitWalk}{-1.30}
\newcommand{\RaoSixtyOneFitFine}{-0.43}
\newcommand{\RaoSixtyOneFitLim}{-0.22}
\newcommand{\RaoSixtyOneFitMinusFitWalk}{+0.00}
\newcommand{\RaoSixtyOneFitMinusFitLim}{+0.00}
% generated by tests/t18_numbers.py from out/t17_axial.json -- do not edit
\newcommand{\AxialPairJ}{0.16}
\newcommand{\AxialPairJwalk}{0.01}
\newcommand{\AxialPairWall}{0.02}
\newcommand{\AxialYard}{9.05}
\newcommand{\AxialBumpWalk}{-10.9}
\newcommand{\AxialBumpFine}{-3.31}
\newcommand{\AxialBumpLim}{-3.66}
\newcommand{\AxialBumpMinusFitWalk}{+0.45}
\newcommand{\AxialBumpMinusFitLim}{-0.86}
\newcommand{\AxialBumpTrials}{4}
\newcommand{\AxialBumpRecords}{4}
\newcommand{\AxialBumpTime}{63\,min}
\newcommand{\AxialBumpStop}{converged}
\newcommand{\AxialBumpPred}{0.0002}
\newcommand{\AxialHalfWalk}{-10.9}
\newcommand{\AxialHalfFine}{-3.47}
\newcommand{\AxialHalfLim}{-3.80}
\newcommand{\AxialHalfMinusFitWalk}{+0.44}
\newcommand{\AxialHalfMinusFitLim}{-1.00}
\newcommand{\AxialHalfTrials}{9}
\newcommand{\AxialHalfRecords}{9}
\newcommand{\AxialHalfTime}{46\,min}
\newcommand{\AxialHalfStop}{converged}
\newcommand{\AxialHalfPred}{0.008}
\newcommand{\AxialFitWalk}{-11.4}
\newcommand{\AxialFitFine}{-2.31}
\newcommand{\AxialFitLim}{-2.80}
\newcommand{\AxialFitMinusFitWalk}{+0.00}
\newcommand{\AxialFitMinusFitLim}{+0.00}
\newcommand{\AxialGenoBound}{-47.0}
% generated by tests/t18_numbers.py from out/t17_hth.json -- do not edit
\newcommand{\HthPairJ}{0.01}
\newcommand{\HthPairJwalk}{0.00}
\newcommand{\HthPairWall}{0.03}
\newcommand{\HthYard}{0.41}
\newcommand{\HthBumpWalk}{-1.27}
\newcommand{\HthBumpFine}{-0.90}
\newcommand{\HthBumpLim}{-0.35}
\newcommand{\HthBumpMinusFitWalk}{+0.03}
\newcommand{\HthBumpMinusFitLim}{-0.01}
\newcommand{\HthBumpTrials}{7}
\newcommand{\HthBumpRecords}{7}
\newcommand{\HthBumpTime}{122\,min}
\newcommand{\HthBumpStop}{converged}
\newcommand{\HthBumpPred}{6e-05}
\newcommand{\HthHalfWalk}{-1.27}
\newcommand{\HthHalfFine}{-0.91}
\newcommand{\HthHalfLim}{-0.37}
\newcommand{\HthHalfMinusFitWalk}{+0.02}
\newcommand{\HthHalfMinusFitLim}{-0.02}
\newcommand{\HthHalfTrials}{8}
\newcommand{\HthHalfRecords}{8}
\newcommand{\HthHalfTime}{127\,min}
\newcommand{\HthHalfStop}{converged}
\newcommand{\HthHalfPred}{0.002}
\newcommand{\HthFitWalk}{-1.30}
\newcommand{\HthFitFine}{-0.89}
\newcommand{\HthFitLim}{-0.33}
\newcommand{\HthFitMinusFitWalk}{+0.00}
\newcommand{\HthFitMinusFitLim}{+0.00}
\newcommand{\HthDoneFlag}{}
\newcommand{\HthLbf}{34\,375\,lbf (-0.03 / -0.03\,lbf vs Rao)}
% generated by tests/t18_numbers.py from out/t17_migdal.json -- do not edit
\newcommand{\MigdalPairJ}{0.01}
\newcommand{\MigdalPairJwalk}{0.00}
\newcommand{\MigdalPairWall}{0.01}
\newcommand{\MigdalYard}{0.00}
\newcommand{\MigdalFreeWalk}{+1.84}
\newcommand{\MigdalFreeFine}{-1.33}
\newcommand{\MigdalFreeLim}{-1.73}
\newcommand{\MigdalFreeMinusFitWalk}{+1.84}
\newcommand{\MigdalFreeMinusFitLim}{-1.73}
\newcommand{\MigdalFreeTrials}{5}
\newcommand{\MigdalFreeRecords}{5}
\newcommand{\MigdalFreeTime}{4\,min}
\newcommand{\MigdalFreeStop}{converged}
\newcommand{\MigdalFreePred}{0.0002}
\newcommand{\MigdalClsWalk}{+1.84}
\newcommand{\MigdalClsFine}{-1.32}
\newcommand{\MigdalClsLim}{-1.72}
\newcommand{\MigdalClsMinusFitWalk}{+1.84}
\newcommand{\MigdalClsMinusFitLim}{-1.72}
\newcommand{\MigdalClsTrials}{14}
\newcommand{\MigdalClsRecords}{14}
\newcommand{\MigdalClsTime}{5\,min}
\newcommand{\MigdalClsStop}{converged}
\newcommand{\MigdalClsPred}{0.001}
\newcommand{\MigdalFitWalk}{+0.00}
\newcommand{\MigdalFitFine}{+0.00}
\newcommand{\MigdalFitLim}{+0.00}
\newcommand{\MigdalFitMinusFitWalk}{+0.00}
\newcommand{\MigdalFitMinusFitLim}{+0.00}
% generated by tests/t18_numbers.py from out/t17_bell.json -- do not edit
\newcommand{\BellPairJ}{0.22}
\newcommand{\BellPairJwalk}{0.00}
\newcommand{\BellPairWall}{0.03}
\newcommand{\BellPairJFine}{0.01}
\newcommand{\BellPairWallFine}{0.02}
\newcommand{\BellYard}{0.00}
\newcommand{\BellTicWalk}{+5.46}
\newcommand{\BellTicFine}{-2.79}
\newcommand{\BellTicLim}{-3.29}
\newcommand{\BellTicMinusFitWalk}{+5.46}
\newcommand{\BellTicMinusFitLim}{-3.29}
\newcommand{\BellTicTrials}{9}
\newcommand{\BellTicRecords}{9}
\newcommand{\BellTicTime}{20\,min}
\newcommand{\BellTicStop}{converged}
\newcommand{\BellTicPred}{0.003}
\newcommand{\BellTicEightOneWalk}{+0.75}
\newcommand{\BellTicEightOneFine}{+0.34}
\newcommand{\BellTicEightOneLim}{+0.26}
\newcommand{\BellTicEightOneMinusFitWalk}{+0.75}
\newcommand{\BellTicEightOneMinusFitLim}{+0.26}
\newcommand{\BellTicEightOneTrials}{7}
\newcommand{\BellTicEightOneRecords}{7}
\newcommand{\BellTicEightOneTime}{44\,min}
\newcommand{\BellTicEightOneStop}{converged}
\newcommand{\BellTicEightOnePred}{0.003}
\newcommand{\BellSteepWalk}{+5.46}
\newcommand{\BellSteepFine}{-3.01}
\newcommand{\BellSteepLim}{-3.52}
\newcommand{\BellSteepMinusFitWalk}{+5.46}
\newcommand{\BellSteepMinusFitLim}{-3.52}
\newcommand{\BellSteepTrials}{10}
\newcommand{\BellSteepRecords}{10}
\newcommand{\BellSteepTime}{17\,min}
\newcommand{\BellSteepStop}{converged}
\newcommand{\BellSteepPred}{0.004}
\newcommand{\BellSteepEightOneWalk}{+0.75}
\newcommand{\BellSteepEightOneFine}{+0.35}
\newcommand{\BellSteepEightOneLim}{+0.27}
\newcommand{\BellSteepEightOneMinusFitWalk}{+0.75}
\newcommand{\BellSteepEightOneMinusFitLim}{+0.27}
\newcommand{\BellSteepEightOneTrials}{10}
\newcommand{\BellSteepEightOneRecords}{10}
\newcommand{\BellSteepEightOneTime}{22\,min}
\newcommand{\BellSteepEightOneStop}{converged}
\newcommand{\BellSteepEightOnePred}{0.004}
\newcommand{\BellFitWalk}{+0.00}
\newcommand{\BellFitFine}{+0.00}
\newcommand{\BellFitLim}{+0.00}
\newcommand{\BellFitMinusFitWalk}{+0.00}
\newcommand{\BellFitMinusFitLim}{+0.00}
\newcommand{\BellThBFit}{16.92}
\newcommand{\BellThBTic}{16.54}
\newcommand{\BellThBTicEightOne}{16.77}
\newcommand{\BellThBSteep}{16.53}
\newcommand{\BellThBSteepEightOne}{16.78}
\newcommand{\BellThBStartTic}{12.88}
\newcommand{\BellThBStartTicEightOne}{12.88}
\newcommand{\BellThBStartSteep}{28.69}
\newcommand{\BellThBStartSteepEightOne}{28.69}
% generated by tests/t18_numbers.py from out/t17_bellideal.json -- do not edit
\newcommand{\BellIdealPairJ}{34.2}
\newcommand{\BellIdealPairJwalk}{15.4}
\newcommand{\BellIdealPairWall}{26.28}
\newcommand{\BellIdealPairJFine}{41.8}
\newcommand{\BellIdealPairWallFine}{28.75}
\newcommand{\BellIdealYard}{0.00}
\newcommand{\BellIdealStretchedWalk}{+24.8}
\newcommand{\BellIdealStretchedFine}{-44.7}
\newcommand{\BellIdealStretchedLim}{-49.3}
\newcommand{\BellIdealStretchedMinusFitWalk}{+24.8}
\newcommand{\BellIdealStretchedMinusFitLim}{-49.3}
\newcommand{\BellIdealStretchedTrials}{13}
\newcommand{\BellIdealStretchedRecords}{14}
\newcommand{\BellIdealStretchedTime}{8\,min}
\newcommand{\BellIdealStretchedStop}{resolution}
\newcommand{\BellIdealStretchedPred}{2}
\newcommand{\BellIdealStretchedEightOneWalk}{-28.8}
\newcommand{\BellIdealStretchedEightOneFine}{-42.0}
\newcommand{\BellIdealStretchedEightOneLim}{-46.4}
\newcommand{\BellIdealStretchedEightOneMinusFitWalk}{-28.8}
\newcommand{\BellIdealStretchedEightOneMinusFitLim}{-46.4}
\newcommand{\BellIdealStretchedEightOneTrials}{12}
\newcommand{\BellIdealStretchedEightOneRecords}{12}
\newcommand{\BellIdealStretchedEightOneTime}{20\,min}
\newcommand{\BellIdealStretchedEightOneStop}{converged}
\newcommand{\BellIdealStretchedEightOnePred}{0.006}
\newcommand{\BellIdealExtendedWalk}{+9.39}
\newcommand{\BellIdealExtendedFine}{-10.5}
\newcommand{\BellIdealExtendedLim}{-12.0}
\newcommand{\BellIdealExtendedMinusFitWalk}{+9.39}
\newcommand{\BellIdealExtendedMinusFitLim}{-12.0}
\newcommand{\BellIdealExtendedTrials}{7}
\newcommand{\BellIdealExtendedRecords}{7}
\newcommand{\BellIdealExtendedTime}{4\,min}
\newcommand{\BellIdealExtendedStop}{converged}
\newcommand{\BellIdealExtendedPred}{0.003}
\newcommand{\BellIdealExtendedEightOneWalk}{+0.77}
\newcommand{\BellIdealExtendedEightOneFine}{-0.20}
\newcommand{\BellIdealExtendedEightOneLim}{-0.58}
\newcommand{\BellIdealExtendedEightOneMinusFitWalk}{+0.77}
\newcommand{\BellIdealExtendedEightOneMinusFitLim}{-0.58}
\newcommand{\BellIdealExtendedEightOneTrials}{12}
\newcommand{\BellIdealExtendedEightOneRecords}{12}
\newcommand{\BellIdealExtendedEightOneTime}{27\,min}
\newcommand{\BellIdealExtendedEightOneStop}{converged}
\newcommand{\BellIdealExtendedEightOnePred}{0.001}
\newcommand{\BellIdealFitWalk}{+0.00}
\newcommand{\BellIdealFitFine}{+0.00}
\newcommand{\BellIdealFitLim}{+0.00}
\newcommand{\BellIdealFitMinusFitWalk}{+0.00}
\newcommand{\BellIdealFitMinusFitLim}{+0.00}
\newcommand{\BellIdealThBFit}{12.62}
\newcommand{\BellIdealThBStretched}{8.15}
\newcommand{\BellIdealThBStretchedEightOne}{8.16}
\newcommand{\BellIdealThBExtended}{12.08}
\newcommand{\BellIdealThBExtendedEightOne}{12.53}
\newcommand{\BellIdealThBStartStretched}{9.14}
\newcommand{\BellIdealThBStartStretchedEightOne}{9.14}
\newcommand{\BellIdealThBStartExtended}{13.49}
\newcommand{\BellIdealThBStartExtendedEightOne}{13.49}

% generated by tests/t18_numbers.py: fallbacks shown in red when a number is missing
\providecommand{\AxialBumpMinusFitLim}{\textcolor{red}{??}}
\providecommand{\AxialBumpMinusFitWalk}{\textcolor{red}{??}}
\providecommand{\AxialBumpTrials}{\textcolor{red}{??}}
\providecommand{\AxialGenoBound}{\textcolor{red}{??}}
\providecommand{\AxialHalfTrials}{\textcolor{red}{??}}
\providecommand{\AxialPairJ}{\textcolor{red}{??}}
\providecommand{\AxialPairWall}{\textcolor{red}{??}}
\providecommand{\BellIdealExtendedEightOneMinusFitLim}{\textcolor{red}{??}}
\providecommand{\BellIdealExtendedEightOneMinusFitWalk}{\textcolor{red}{??}}
\providecommand{\BellIdealPairJFine}{\textcolor{red}{??}}
\providecommand{\BellIdealStretchedEightOneMinusFitLim}{\textcolor{red}{??}}
\providecommand{\BellIdealThBFit}{\textcolor{red}{??}}
\providecommand{\BellIdealThBStretchedEightOne}{\textcolor{red}{??}}
\providecommand{\BellPairJFine}{\textcolor{red}{??}}
\providecommand{\BellPairWallFine}{\textcolor{red}{??}}
\providecommand{\BellThBFit}{\textcolor{red}{??}}
\providecommand{\BellThBStartSteep}{\textcolor{red}{??}}
\providecommand{\BellThBStartTic}{\textcolor{red}{??}}
\providecommand{\BellThBSteepEightOne}{\textcolor{red}{??}}
\providecommand{\BellThBTicEightOne}{\textcolor{red}{??}}
\providecommand{\BellTicEightOneMinusFitLim}{\textcolor{red}{??}}
\providecommand{\BellTicMinusFitLim}{\textcolor{red}{??}}
\providecommand{\BellTicMinusFitWalk}{\textcolor{red}{??}}
\providecommand{\HthBumpFine}{\textcolor{red}{??}}
\providecommand{\HthBumpMinusFitLim}{\textcolor{red}{??}}
\providecommand{\HthHalfFine}{\textcolor{red}{??}}
\providecommand{\HthLbf}{\textcolor{red}{??}}
\providecommand{\HthPairWall}{\textcolor{red}{??}}
\providecommand{\HthProv}{\textcolor{red}{??}}
\providecommand{\MethaneAMinusFitLim}{\textcolor{red}{??}}
\providecommand{\MethaneAMinusFitWalk}{\textcolor{red}{??}}
\providecommand{\MethaneATrials}{\textcolor{red}{??}}
\providecommand{\MethaneCTrials}{\textcolor{red}{??}}
\providecommand{\MethanePairJ}{\textcolor{red}{??}}
\providecommand{\MethanePairWall}{\textcolor{red}{??}}
\providecommand{\MigdalClsTrials}{\textcolor{red}{??}}
\providecommand{\MigdalFreeFine}{\textcolor{red}{??}}
\providecommand{\MigdalFreeMinusFitLim}{\textcolor{red}{??}}
\providecommand{\MigdalFreeTrials}{\textcolor{red}{??}}
\providecommand{\MigdalFreeWalk}{\textcolor{red}{??}}
\providecommand{\MigdalPairJ}{\textcolor{red}{??}}
\providecommand{\MigdalPairWall}{\textcolor{red}{??}}
\providecommand{\PlanarBFine}{\textcolor{red}{??}}
\providecommand{\PlanarBMinusFitLim}{\textcolor{red}{??}}
\providecommand{\PlanarFitFine}{\textcolor{red}{??}}
\providecommand{\PlanarPairJ}{\textcolor{red}{??}}
\providecommand{\PlanarPairWall}{\textcolor{red}{??}}
\providecommand{\RaoSixtyOneBumpMinusFitLim}{\textcolor{red}{??}}
\providecommand{\RaoSixtyOneBumpMinusFitWalk}{\textcolor{red}{??}}
\providecommand{\RaoSixtyOneBumpTime}{\textcolor{red}{??}}
\providecommand{\RaoSixtyOneBumpTrials}{\textcolor{red}{??}}
\providecommand{\RaoSixtyOneHalfTime}{\textcolor{red}{??}}
\providecommand{\RaoSixtyOneHalfTrials}{\textcolor{red}{??}}
\providecommand{\RaoSixtyOneOldBumpMinusFitWalk}{\textcolor{red}{??}}
\providecommand{\RaoSixtyOneOldHalfMinusFitWalk}{\textcolor{red}{??}}
\providecommand{\RaoSixtyOnePairJ}{\textcolor{red}{??}}
\providecommand{\RaoSixtyOnePairWall}{\textcolor{red}{??}}

% verdicts and plain-words explanations of the slide deck (hand-written; every number in them is a macro generated by
% tests/t18_numbers.py from the measured JSONs, or is quoted from a log named next to it)
\newcommand{\PlanarFound}{\yes}
\newcommand{\MethaneFound}{\yes}
\newcommand{\RaoSixtyOneFound}{\yes}
\newcommand{\AxialFound}{\yes}
\newcommand{\BellFound}{\yes}
\newcommand{\BellIdealFound}{\partly}
\newcommand{\MigdalFound}{\yes}
\newcommand{\PLANARVERDICT}{Four starts --- two distant ones, and two with the no-fold rule on (one starting outside it)
--- land on one value (\PlanarPairJ\ppm\ apart) and one wall (\PlanarPairWall\,mm per m), in 5--13 steps, 2--4\,min.
At the finest grid the result is \PlanarBFine\ppm\ from the bound, the closest 12-point curve \PlanarFitFine\ppm; in the
limit the result sits \PlanarBMinusFitLim\ppm\ from that curve. (Old driver: 27--28 segments with a radius floor 16
times finer than its default --- with the default it stopped 1\ppm\ short; with the rule on it stalled at $-49$\ppm, and
from a start breaking the rule it did not move.)}
\newcommand{\PLANARWHY}{The exact wall bends abruptly where the first lip wave lands. A smooth curve through 12
points cannot follow that bend: its wall angle is off by $\delta\theta\approx2\times10^{-3}$ rad on average, and a wall
angle error costs about $2\,\delta\theta^2\approx9$\ppm\ of thrust. No optimiser can beat the curve it is given; the
walls of the result and of the exact plug differ by $\sim$1\,mm per m where the thrust does not care (flat top).}
\newcommand{\METHANEVERDICT}{\yes. Two starts land on one point (\MethanePairJ\ppm,
walls \MethanePairWall\,mm per m apart) in \MethaneATrials\ and \MethaneCTrials\ steps: \MethaneAMinusFitWalk\ppm\ from
the closest curve of Rao's contour on the optimiser's grid, \MethaneAMinusFitLim\ppm\ in the limit. Same landing as the old
driver; this is the RDE-like plug ($\theta_\mathrm{wall}+\mu\ge5.5^\circ$, far from the steep regime of Rao 1961).}
\newcommand{\METHANEWHY}{Rao's method gives the best plug for exactly this inflow,
length and end point; the optimiser, which knows none of that, finds the same wall to a fraction of a millimetre per
metre. Above Rao on the coarse grid, below it on finer grids: the grid effect of every case, not a gain.}
\newcommand{\RAOVERDICT}{\yes, with the default settings. Both starts land on one point
(\RaoSixtyOnePairJ\ppm, walls \RaoSixtyOnePairWall\,mm per m apart) in \RaoSixtyOneBumpTrials\ and
\RaoSixtyOneHalfTrials\ steps: \RaoSixtyOneBumpMinusFitWalk\ppm\ from the closest curve of Rao's contour on the optimiser's grid,
\RaoSixtyOneBumpMinusFitLim\ppm\ in the limit; the wall is 0.4\,mm per m from Rao's. With the old driver the same starts
stopped \RaoSixtyOneOldBumpMinusFitWalk\ and \RaoSixtyOneOldHalfMinusFitWalk\ppm\ from that curve; it got there only with
hand-set short steps and two more hours per walk.}
\newcommand{\RAOWHY}{Rao's $\varepsilon=3.81$ plug is so steep near the throat that the
characteristics at the wall point slightly downward, and the thrust varies along a narrow ridge. The old driver walked
across the ridge and kept falling off it; the new one measures the ridge's shape once, keeps it, and walks along it.}
\newcommand{\AXIALVERDICT}{\yes. Two starts land on one point (\AxialPairJ\ppm,
walls \AxialPairWall\,mm per m apart) in \AxialBumpTrials\ and \AxialHalfTrials\ steps: \AxialBumpMinusFitWalk\ppm\ from
the closest curve of Rao's contour on the optimiser's grid, \AxialBumpMinusFitLim\ppm\ in the limit.}
\newcommand{\AXIALWHY}{The flow leaves this plug uniform and axial at ambient pressure: the most thrust any plug
can give. Rao's contour delivers it (to the consistency of GENO's data) and the optimiser finds Rao's contour; the
small remaining gap is the 9-point curve's and the grid's.}
\newcommand{\BELLVERDICT}{\yes. Two starts on opposite sides of the optimum (attachment \BellThBStartTic\textdegree\ and
\BellThBStartSteep\textdegree, thrust $-2.2\times10^{-3}$ and $-3.1\times10^{-2}$) end on one wall and one value. Optimised
on the grid twice as fine: attachment \BellThBTicEightOne\textdegree\ and \BellThBSteepEightOne\textdegree\ (closest curve
to GENO's \BellThBFit\textdegree), walls \BellPairWallFine\,mm per m apart, \BellPairJFine\ppm\ apart; in the limit
\BellTicEightOneMinusFitLim\ppm\ from the closest curve. Optimised on the coarser grid the same starts end
\BellTicMinusFitWalk\ppm\ above that curve on their own grid but \BellTicMinusFitLim\ppm\ in the limit: that grid is too
coarse to rank walls at the ppm level.}
\newcommand{\BELLWHY}{Rao's bell is the optimum; ours --- a curve through 8 points --- follows GENO's within
0.6\,mm per m. The optimiser sees the thrust only through its grid, and that grid's small error changes with the wall by a
few ppm: it gets used. One grid halving divides that effect by about 10 (from $-3.4$ to $+0.3$\ppm\ here).}
\newcommand{\BELLIDEALVERDICT}{\partly. From Rao's contour for length 6 followed by a cylinder the optimiser finds the
ideal bell: on the grid twice as fine \BellIdealExtendedEightOneMinusFitWalk\ppm\ from its closest curve on its own grid,
\BellIdealExtendedEightOneMinusFitLim\ppm\ in the limit. From the same contour stretched to the full length (attachment
9.1\textdegree) it converges --- correctly, its curvature measured there --- to a \textbf{second top}: attachment
\BellIdealThBStretchedEightOne\textdegree\ instead of \BellIdealThBFit\textdegree, wall 3\,\% of the throat radius away,
\BellIdealStretchedEightOneMinusFitLim\ppm\ in the limit. The two-start check flags it (\BellIdealPairJFine\ppm\ apart).}
\newcommand{\BELLIDEALWHY}{At its full length the ideal bell's thrust is extremely flat: walls up to 3\,\% of the throat
radius apart differ by only about 45\ppm. Both results are strict local maxima (all curvatures negative, every probing march lower) and between
them the thrust dips by 14\ppm\ (backup slide): two real tops of the computed thrust; a local optimiser stops on the one
it climbs. Not the 8-point wall (with 12 points the same start
still does not get there), not the throat-arc panel the objective leaves out, not the axial-Mach margin; the cause is
open. Two distant starts are the protection, and here they flag it.}
\newcommand{\MIGDALVERDICT}{\yes. With and without the no-fold rule --- the walk with the rule starting \emph{outside} it --- the two walks land on one point (\MigdalPairJ\ppm, walls \MigdalPairWall\,mm per m apart) in \MigdalFreeTrials\ and
\MigdalClsTrials\ steps. On the optimiser's grid \MigdalFreeWalk\ppm\ from GENO's perfect nozzle, on the finest grid
\MigdalFreeFine\ppm\ (limit \MigdalFreeMinusFitLim\ppm). The production walker (Newton metric, no-fold rule on, its own radius
rule) lands on the same point, $C_F$ equal to 9 digits, in 3 steps. Old driver: from a start breaking the rule it did not move.}
\newcommand{\MIGDALWHY}{The perfect nozzle turns all the flow axial at the exit; walls that do the same within the
grid's error give the same thrust. The result and GENO's walls differ by 0.5\,mm per m on the plug and 0.1 on the shroud,
where the thrust cannot tell them apart; the ppm left are the grid's, as in every case.}
\newcommand{\OPENONE}{\textbf{The grid sets the precision.} The walk maximises the thrust its grid computes and
uses that grid's small error: 1--2\ppm\ on the plug grids, $\sim$3--5\ppm\ on the bell engine at $41\times120$, more on the
flat ideal bell. One grid halving cuts the effect by about 10 (bell TOC: $-3.4\to+0.3$\ppm). Rule: walk on a grid whose
error varies with the wall less than the precision wanted, and re-march the result once on a twice finer grid. When the
grid's thrust is not smooth at the scale of the last steps the driver now says so and stops (``RESOLUTION''); that
stop was added when the coarse ideal-bell walk climbed grid steps, and replayed on the logs of every other walk here it
never fires.}
\newcommand{\OPENTWO}{\textbf{Two tops.} A local optimiser stops on the top it climbs: the ideal bell at its full
length has a second top 30--45\ppm\ lower (a 14\ppm\ valley between the two results); its cause is not identified (not the
wall description, not the throat panel, not the axial-Mach margin). Protection: always two distant starts; a third when
they disagree. Every other case here has one top (all starts agree to $<0.3$\ppm).}
\newcommand{\OPENTHREE}{\textbf{Starts and cost.} The bell engine refuses a conical start (internal shock,
outside the model) and long shallow starts (a cell's axial Mach below 1 near the throat). The walker's generic start on the
annular nozzle folds before the driver is called (the walker's own start logic, unchanged); the driver itself now restores
starts that break the no-fold rule. \textbf{Cost:} on the Rao posings one gradient takes 40--100\,s (today's s2 was loaded:
our runs got 25--35\,\% of a core), and the two measured curvatures (2 $\times$ 18 gradients) are $\sim$90\,\% of a walk:
1--1.7\,h today. One-sided differences would halve that, a compiled replay would cut it further; not done, so that the
runs shown here match the code that is now the default. Other cases: 2--45\,min per walk.}
\newcommand{\BELLVSAH}{$+0.3$\ppm\ (optimised on the grid twice as fine) to $-3.5$\ppm\ (coarser grid)}
\newcommand{\HthFound}{\ifdefined\HthDoneFlag\yes\else\yes$^*$\fi}
\newcommand{\HTHVERDICTFINAL}{\yes. From a characteristic of Rao's own field, Humphreys' method gives a contour ``almost
identical to Rao's'' ($-2$\,lbf); our optimiser, from two starts, lands \HthBumpFine\ and \HthHalfFine\ppm\ from Rao's
contour on the finest grid (\HthBumpMinusFitLim\ppm\ from its closest curve in the limit), walls \HthPairWall\,mm per m
apart.}
\newcommand{\HTHVERDICTPROV}{\yes\ on the optimiser's grid. Their Table 3 lies on GENO's Rao contour within
$5\times10^{-5}R_E$ on the designed part; Humphreys' method from a characteristic of Rao's field: $-2$\,lbf
($-58$\ppm). Our SQP on the same problem: \HthProv. The check on finer grids is still running.}
\newcommand{\HTHVERDICT}{\ifdefined\HthDoneFlag\HTHVERDICTFINAL\else\HTHVERDICTPROV\fi}
\newcommand{\HTHWHY}{Same inflow, length and end point: one best plug. Rao builds it from the exit characteristic,
Humphreys from the given start line, our optimiser by climbing on thrust; all three find it. Their design 4\,\% lower
is for their own engine (another start line, wall friction): a different problem.}


\title{Does the optimiser find the optimum?}
\subtitle{rde-nozzle revalidation: plugs, bells, annular nozzle}
\author{harness \texttt{RDE\_Design/validation/\_revalidation\_2026-10-07/}}
\date{2026-10-08}

\begin{document}

% =====================================================================================================================
\begin{frame}
\titlepage
\vspace{-6mm}
\begin{center}\small
Every case has an optimum known \emph{without} the model: an exact bound, a paper, or Rao's method in GENO
(branch \texttt{exp/ws-c-prod-nozzles}, validated).\\[1mm]
The optimiser is run from two distant starts with \textbf{one setting for every case}; it is never told the optimum.
\end{center}
\end{frame}

% =====================================================================================================================
\begin{frame}{The answer, case by case}
\centering\small
\resizebox{0.97\linewidth}{!}{% generated by tests/t18_numbers.py -- do not edit
\begin{tabular}{@{}llrrc@{}}\toprule
case & the known optimum & result vs optimum, two starts (finest grid) & starts differ by & found?\\\midrule
Planar plug & exact ideal plug, $T=\dot mV_a$ & -10.3 / -10.3\ppm\ \ (closest curve -8.96) & 0.01\ppm & \PlanarFound\\
Rao plug, methane (RDE-like) & Rao (GENO) & -2.11 / -2.18\ppm\ \ (closest curve -0.24) & 0.07\ppm & \MethaneFound\\
Rao 1961 plug, $\varepsilon$ 3.81 & Rao 1961 Table 1 (= GENO) & -1.30 / -1.35\ppm\ \ (closest curve -0.43) & 0.05\ppm & \RaoSixtyOneFound\\
Axial-exit plug & Rao $\theta_E=0$, $T=\dot mV_a$ & -3.31 / -3.47\ppm\ \ (closest curve -2.31) & 0.16\ppm & \AxialFound\\
Humphreys 1971 plug & Rao (their Table 3 = GENO) & -0.90 / -0.91\ppm\ \ (closest curve -0.89) & 0.01\ppm & \HthFound\\
Bell, length $4\,r_t$ (grid twice as fine) & Rao 1958 TOC (GENO) & +0.34 / +0.35\ppm\ \ (closest curve +0.00) & 0.01\ppm & \BellFound\\
Bell, ideal length (grid twice as fine) & ideal bell, $F_\mathrm{1D}$ & -42.0 / -0.20\ppm\ \ (closest curve +0.00) & 41.8\ppm & \BellIdealFound\\
Annular nozzle & Migdal (GENO), $F_\mathrm{1D}$ & -1.33 / -1.32\ppm\ \ (closest curve +0.00) & 0.01\ppm & \MigdalFound\\
\bottomrule
\end{tabular}
}

\vspace{3mm}
\begin{minipage}{0.95\linewidth}\footnotesize
\textbf{Result vs optimum}: the optimiser's thrust minus the optimum's, re-computed on the finest of three grids, for each
of the two starts, in \textbf{ppm} (parts per million of thrust: 1\ppm\ = 1\,N per MN). The optimiser describes a wall
by a smooth curve through 8--12 points; the closest such curve to the optimum is itself a little below a sharp optimum
(in brackets): that is about the most any walk can reach. Bells: measured against the closest curve to GENO's wall (the
bell engine marches only its own curves). \textbf{Found?}: both starts reach the optimum within the grid and the curve
description; ``partly'': one start does, the other stops on a second, lower top. \textbf{Times} in the case tables: s2
was shared with other users' jobs at normal priority from $\sim$13:20; our runs got 25--35\,\% of a core.
\end{minipage}
\end{frame}

% =====================================================================================================================
\begin{frame}{The optimiser's checks, case by case}
\centering
\resizebox{0.86\linewidth}{!}{% generated by tests/t18_numbers.py -- do not edit
\begin{tabular}{@{}llcccc@{}}\toprule
case & start & C1 converged & C2 reaches its curve & C3 starts agree & C4 not above optimum\\\midrule
Planar plug & 15\textdegree\ fan plug & \textcolor{green!45!black}{\textbf{pass}} {\tiny converged} & \textcolor{green!45!black}{\textbf{pass}} {\tiny +1.70} & \textcolor{green!45!black}{\textbf{pass}} {\tiny 0.01$\le$2.69} & \textcolor{green!45!black}{\textbf{pass}} {\tiny -10.3$\le$1.91}\\
 & straight line & \textcolor{green!45!black}{\textbf{pass}} {\tiny converged} & \textcolor{green!45!black}{\textbf{pass}} {\tiny +1.71} & \textcolor{green!45!black}{\textbf{pass}} {\tiny 0.01$\le$2.69} & \textcolor{green!45!black}{\textbf{pass}} {\tiny -10.3$\le$1.91}\\
Methane plug & deformed Rao & \textcolor{green!45!black}{\textbf{pass}} {\tiny converged} & \textcolor{green!45!black}{\textbf{pass}} {\tiny +1.95} & \textcolor{green!45!black}{\textbf{pass}} {\tiny 0.07$\le$3.83} & \textcolor{green!45!black}{\textbf{pass}} {\tiny -2.11$\le$0.06}\\
 & straight line & \textcolor{green!45!black}{\textbf{pass}} {\tiny converged} & \textcolor{green!45!black}{\textbf{pass}} {\tiny +1.95} & \textcolor{green!45!black}{\textbf{pass}} {\tiny 0.07$\le$3.83} & \textcolor{green!45!black}{\textbf{pass}} {\tiny -2.18$\le$0.06}\\
Rao 1961 & deformed Rao & \textcolor{green!45!black}{\textbf{pass}} {\tiny converged} & \textcolor{green!45!black}{\textbf{pass}} {\tiny +0.77} & \textcolor{green!45!black}{\textbf{pass}} {\tiny 0.05$\le$1.75} & \textcolor{green!45!black}{\textbf{pass}} {\tiny -1.30$\le$0.26}\\
 & Rao--line average & \textcolor{green!45!black}{\textbf{pass}} {\tiny converged} & \textcolor{green!45!black}{\textbf{pass}} {\tiny +0.77} & \textcolor{green!45!black}{\textbf{pass}} {\tiny 0.05$\le$1.75} & \textcolor{green!45!black}{\textbf{pass}} {\tiny -1.35$\le$0.26}\\
Axial exit & deformed Rao & \textcolor{green!45!black}{\textbf{pass}} {\tiny converged} & \textcolor{green!45!black}{\textbf{pass}} {\tiny +0.45} & \textcolor{green!45!black}{\textbf{pass}} {\tiny 0.16$\le$1.31} & \textcolor{green!45!black}{\textbf{pass}} {\tiny -3.31$\le$0.49}\\
 & Rao--line average & \textcolor{green!45!black}{\textbf{pass}} {\tiny converged} & \textcolor{green!45!black}{\textbf{pass}} {\tiny +0.44} & \textcolor{green!45!black}{\textbf{pass}} {\tiny 0.16$\le$1.31} & \textcolor{green!45!black}{\textbf{pass}} {\tiny -3.47$\le$0.49}\\
Humphreys 1971 & deformed Rao & \textcolor{green!45!black}{\textbf{pass}} {\tiny converged} & \textcolor{green!45!black}{\textbf{pass}} {\tiny +0.03} & \textcolor{green!45!black}{\textbf{pass}} {\tiny 0.01$\le$0.04} & \textcolor{green!45!black}{\textbf{pass}} {\tiny -0.90$\le$0.23}\\
 & Rao--line average & \textcolor{green!45!black}{\textbf{pass}} {\tiny converged} & \textcolor{green!45!black}{\textbf{pass}} {\tiny +0.02} & \textcolor{green!45!black}{\textbf{pass}} {\tiny 0.01$\le$0.04} & \textcolor{green!45!black}{\textbf{pass}} {\tiny -0.91$\le$0.23}\\
Bell TOC & truncated ideal & \textcolor{green!45!black}{\textbf{pass}} {\tiny converged} & \textcolor{green!45!black}{\textbf{pass}} {\tiny +0.75} & \textcolor{green!45!black}{\textbf{pass}} {\tiny 0.01$\le$0.48} & n/a\\
 & squeezed TOC & \textcolor{green!45!black}{\textbf{pass}} {\tiny converged} & \textcolor{green!45!black}{\textbf{pass}} {\tiny +0.75} & \textcolor{green!45!black}{\textbf{pass}} {\tiny 0.01$\le$0.48} & n/a\\
Ideal bell & stretched TOC & \textcolor{green!45!black}{\textbf{pass}} {\tiny converged} & \textcolor{red!70!black}{\textbf{FAIL}} {\tiny -28.8} & \textcolor{red!70!black}{\textbf{FAIL}} {\tiny 41.8$\le$1.35} & n/a\\
 & TOC + cylinder & \textcolor{green!45!black}{\textbf{pass}} {\tiny converged} & \textcolor{green!45!black}{\textbf{pass}} {\tiny +0.77} & \textcolor{red!70!black}{\textbf{FAIL}} {\tiny 41.8$\le$1.35} & n/a\\
Annular & no rule & \textcolor{green!45!black}{\textbf{pass}} {\tiny converged} & \textcolor{green!45!black}{\textbf{pass}} {\tiny +1.84} & \textcolor{green!45!black}{\textbf{pass}} {\tiny 0.01$\le$3.57} & \textcolor{green!45!black}{\textbf{pass}} {\tiny -1.33$\le$0.00}\\
 & no-fold rule on & \textcolor{green!45!black}{\textbf{pass}} {\tiny converged} & \textcolor{green!45!black}{\textbf{pass}} {\tiny +1.84} & \textcolor{green!45!black}{\textbf{pass}} {\tiny 0.01$\le$3.57} & \textcolor{green!45!black}{\textbf{pass}} {\tiny -1.32$\le$0.00}\\
\midrule \multicolumn{6}{r}{57 / 60 checks pass}\\
\bottomrule
\end{tabular}
}

\vspace{1.5mm}
\begin{minipage}{0.97\linewidth}\scriptsize
\textbf{C1} the driver declares convergence (gain still predicted, curvature re-measured, $<10^{-8}$ of thrust).
\textbf{C2} on its own grid the result reaches the closest curve to the optimum (tolerance: the gain still predicted).
\textbf{C3} the two starts agree within the grid's resolution of wall comparisons in that case (the smaller of the
two starts' shifts from its own grid to the limit). \textbf{C4} on the finest grid the result is not above the optimum
beyond what the grid resolves there (shift of the closest curve's distance from the optimum between the two finest grids;
n/a for bells: the bell engine cannot march GENO's wall). Grey: measured value $\le$ tolerance, ppm. The three FAILs are
the ideal bell's second top: a real top of the posed problem (backup slide), not a wrong stop of the driver.
\end{minipage}
\end{frame}

% =====================================================================================================================
\begin{frame}{How to read the numbers}
\begin{columns}[T]
\begin{column}{0.5\textwidth}
\textbf{Three sources of small differences, none of them the optimiser's:}
\begin{itemize}\footnotesize
\item \textbf{Grid.} Thrust is computed on a grid (like measuring with a ruler); halving the spacing divides the error
by about 4. Every result is re-computed on 3 grids and extrapolated.
\item \textbf{Wall description.} 8--12 points + a smooth curve cannot follow a wall whose curvature jumps (where the
first lip wave lands, where the bell leaves the throat arc). Cost: a few ppm.
\item \textbf{Tuning to the ruler.} The walk maximises the thrust \emph{on its grid}, including that grid's error:
\[J_h=J+e_h\quad\Rightarrow\quad \nabla J=-\nabla e_h\ \text{at the landing,}\]
so it can sit \emph{above} the closest curve on the optimiser's grid and 1--2\ppm\ below it on a finer one.
\end{itemize}
\end{column}
\begin{column}{0.47\textwidth}
\textbf{What counts as ``found''}
\begin{enumerate}\footnotesize
\item On the optimiser's own grid the landing is at least as good as the closest curve of the optimum (it did its job on
the objective it sees).
\item Two distant starts land on one point (thrust and wall).
\item On finer grids the landing is within the grid / wall-description effects above of the optimum --- never
\emph{above} the optimum beyond its error (Rao is the optimum of the same problem: a gain would be a defect of ours).
\end{enumerate}
\vspace{1mm}
\plain{Flat top.}{Near an optimum many walls give almost the same thrust: walls 0.3--0.5\,mm per metre apart can differ
by less than 1\ppm. Shapes are therefore compared to that scale, not to zero.}
\plain{No-fold rule.}{The model is shock-free: characteristics of one family must not cross. Where the optimiser could
gain by folding them (a shock in reality), the rule forbids it; it is a constraint of the walk.}
\end{column}
\end{columns}
\end{frame}

% =====================================================================================================================
\begin{frame}{The driver: what was wrong, what changed}
\begin{columns}[T]
\begin{column}{0.49\textwidth}
\textbf{Before} (segmented \texttt{trust-constr}, \texttt{PSPL\_DRIVER=legacy})\\[1mm]
\footnotesize
\begin{enumerate}
\item A segment walks a \emph{frozen} copy of the march, exact only within $\sim10^{-3}$ of its base, but its step is
not bounded by its radius (2.8 at radius $10^{-3}$, measured): it predicted $+1$--$4\,\%$, the real march said
``worse''.
\item The curvature was rebuilt from nothing at every segment: gradient-like steps, slow on a narrow ridge.
\item Stop rule ``two short steps'' read slow progress as the top.
\item Restoration inert (a start breaking the no-fold rule never moved); with the rule on, interior-point segments
returned worse points.
\end{enumerate}
\end{column}
\begin{column}{0.49\textwidth}
\textbf{Now} (\texttt{a1\_qntr}, default)\\[1mm]
\footnotesize
\begin{enumerate}
\item Every trial point is the \textbf{true march}; the frozen copy is used only for the gradient \emph{at} a marched
point, where it is exact.
\item Step from a quadratic model whose curvature $B$ is measured once (secant Hessian) and \textbf{kept} (BFGS):
\[d=\arg\max_{\|d\|\le\Delta}\ g^\top d-\tfrac12 d^\top Bd,\quad \rho=\tfrac{\text{actual gain}}{\text{predicted gain}}.\]
\item Stop only when the gain still predicted, $\tfrac12 g^\top B^{-1}g$, is below $10^{-8}J$ \textbf{and} stays below
with $B$ re-measured there.
\item Restoration = the same machinery on the rule's margin; the rule is a linear constraint of the model.
\end{enumerate}
\end{column}
\end{columns}
\vspace{1mm}
\plain{In plain words.}{Climbing a hill in fog: the old method trusted a map valid for one step but walked ten,
forgot the shape of the hill after every stretch, and stopped when its steps got short. The new one checks every step
on the real ground, remembers the shape of the hill, and stops only when even its best estimate of the remaining
climb is below one part in $10^8$.}
\end{frame}

% =====================================================================================================================
\begin{frame}{Rao 1961: the case that exposed it}
\begin{columns}[T]
\begin{column}{0.56\textwidth}
\includegraphics[width=\linewidth]{../figs/19_rao61_history.png}
\end{column}
\begin{column}{0.42\textwidth}
\footnotesize
Rao's $\varepsilon=3.81$ plug is very steep near the throat: the right-running characteristics at the wall turn
\emph{downwards} ($\theta_\mathrm{wall}+\mu=-3.3^\circ$, every other plug here $>0$). The thrust then varies along a
narrow ridge: the best step along the gradient is $2.4\times10^{-4}$ for a gain of 0.6\ppm.

\medskip
\textbf{Old driver:} ended \RaoSixtyOneOldBumpMinusFitWalk\ and \RaoSixtyOneOldHalfMinusFitWalk\ppm\ from the closest
curve (optimiser's grid). \textbf{New driver, same settings as every case:} \RaoSixtyOneBumpTrials{} and
\RaoSixtyOneHalfTrials{} steps, \RaoSixtyOneBumpTime{} and \RaoSixtyOneHalfTime{} (on a loaded machine), both above that
curve on their grid.

\medskip
\why{The old method walked along a knife-edge ridge with steps that kept falling off it; after each fall it took
shorter steps until they looked like ``no progress''. Remembering the ridge's shape (the curvature) lets the new method
step along the ridge instead of across it.}
\end{column}
\end{columns}
\end{frame}

% =====================================================================================================================
\begin{frame}{What the optimiser stands on (model checks)}
\centering\footnotesize
\resizebox{\linewidth}{!}{\begin{tabular}{@{}llll@{}}\toprule
check & against & result & \\\midrule
gas tables & exact & $p,\rho$ within $1.2\times10^{-7}$ (bound $4.9\times10^{-7}$) & \yes\\
one characteristic cell & two exact fields & local order 3.0 & \yes\\
planar plug, centred expansion & exact & order 2.0; thrust $-0.28$\ppm\ at the finest grid & \yes\\
axisymmetric two-wall channel & exact source flow & order 2.00; thrust within 10\ppm & \yes\\
plug + shroud on GENO's Migdal walls & GENO, $F_\mathrm{1D}$ & wall $p$ order 2; thrust limit $-1.6$\ppm\ from $F_\mathrm{1D}$ & \yes\\
plug on GENO's Rao (methane) & GENO & wall $p$ to $7\times10^{-5}$ = GENO's own consistency & \yes\\
\textbf{bell} march on GENO's TOC and ideal bell & grids, GENO & order 1.8 / 2.4; $\Delta C_F$(TOC$-$ideal) $=-0.008484$ vs GENO $-0.0084$ & \yes\\
gradients (reverse AD) & finite differences & equal within $10^{-5}$; replay $=$ march to $10^{-16}$ & \yes\\
Rao 1961, Humphreys 1971 & the papers & tables reproduced within printing (18/20; both misses the paper's) & \yes\\\bottomrule
\end{tabular}}

\vspace{2mm}
\begin{minipage}{0.9\linewidth}\footnotesize
The march is the method of characteristics: second-order accurate, exact solutions recovered at order 2, GENO's walls
and pressures recovered to GENO's own consistency. Where a cell is nearly singular (right-running characteristic
horizontal at the wall) its Newton solve could stop one step short; the rescue that fixes it is now \textbf{on by
default} (certified cells unchanged bit for bit). Re-run on 2026-10-09 with the committed code; every check now carries a \emph{rejector} (a deliberately corrupted model must fail it): all 10 caught. Details: \texttt{doc/validation.pdf} (first day, partly superseded) and \texttt{doc/TRIAGE\_2026-10-09.md}.
\end{minipage}
\end{frame}

% =====================================================================================================================
\caseframe{Planar plug vs the exact ideal plug}{17_planar.png}{With the jet edge and the base at $p_a$, per unit mass
$u+\frac{p-p_a}{\rho u}\le V_a$ (supersonic flow), so $T\le\dot mV_a$, with equality only for a uniform axial exit:
the exact ideal plug.}{num_planar_table}{\yes. \PLANARVERDICT}{\PLANARWHY}

% =====================================================================================================================
\caseframe{Rao plug, methane gas, lip Mach 1.556 (RDE-like)}{17_methane.png}{GENO's Rao contour: the model starts on GENO's last internal characteristic and designs the
plug from $x_s$ (where the first lip wave lands) to the end $D$; Rao's contour is the maximum of exactly this problem.}{num_methane_table}{\METHANEVERDICT}{\METHANEWHY}

% =====================================================================================================================
\caseframe{Rao 1961, $\varepsilon=3.81$ (the paper's Table 1)}{17_rao61.png}{Rao's contour (his Table 1 = GENO's, within printing); same posing as the methane plug,
$\gamma=1.23$, $M_E=2.4$, $\theta_E=-8.25^\circ$.}{num_rao61_table}{\RAOVERDICT}{\RAOWHY}

% =====================================================================================================================
\caseframe{Axial exit: Rao's $\theta_E=0$ member}{17_axial.png}{With $\theta_E=0$ Rao's conditions give $C_3=0$, so $\theta\equiv0$ on the control surface:
uniform axial flow at $p_a$, the bound $T=\dot mV_a$. GENO cannot reach the axis, so the plug ends at $R_D=0.4$ (base
at $p_a$). GENO's contour reaches the bound to \AxialGenoBound\ppm\ in the limit: the start line is GENO's own data, which
close mass and momentum against GENO's exit surface only to $10^{-5}$--$10^{-4}$ (measured on the methane case);
results are compared with GENO's contour on that same line.}{num_axial_table}{\AXIALVERDICT}{\AXIALWHY}

% =====================================================================================================================
\begin{frame}{Rao vs Humphreys et al.\ 1971 vs our optimiser: the same plug}
\begin{columns}[T]
\begin{column}{0.53\textwidth}
\includegraphics[width=\linewidth,height=0.47\textheight,keepaspectratio]{../figs/17_hth.png}\\[0.5mm]
\resizebox{\linewidth}{!}{\begin{tabular}{@{}llr@{}}\toprule
method & start line & thrust\\\midrule
Rao's method (their Table 3) & Rao's sonic line & 34\,253\,lbf\\
\quad the same, our solver of Rao's equations & & 34\,254\,lbf\\\midrule
Humphreys' variational method & a characteristic of Rao's field & 34\,373\,lbf\\
\quad ``compared to \dots\ for Rao's nozzle'' (their words) & & 34\,375\,lbf\\
our optimiser (2 starts): ratio to Rao's contour $\times$ 34\,375 & GENO's characteristic of Rao's field & \HthLbf\\\midrule
Humphreys' fixed-inlet optimum & their own start line, wall friction & 32\,881\,lbf\\\bottomrule
\end{tabular}}\\[0.5mm]
\resizebox{\linewidth}{!}{% generated by tests/t18_numbers.py from out/t17_hth.json -- do not edit
\begin{tabular}{@{}lrrrrrr@{}}\toprule
 & \multicolumn{2}{c}{vs the optimum} & \multicolumn{2}{c}{vs its closest curve} & & \\
start & optimiser's grid & finest & optimiser's grid & limit & steps & time\\\midrule
Rao's contour deformed (up to 1\,\%) & -1.27 & -0.90 & +0.03 & -0.01$\pm$0.01 & 7 & 122\,min\\
average of Rao's contour and a straight line & -1.27 & -0.91 & +0.02 & -0.02$\pm$0.01 & 8 & 127\,min\\
\midrule closest curve & -1.30 & -0.89 & 0 & 0 & & \\
\bottomrule
\end{tabular}
}
\end{column}
\begin{column}{0.45\textwidth}
\scriptsize
\textbf{Their Rao nozzle:} 148.08\,lbm/s, 500\,psia, 6000\,$^\circ$R, $\gamma$ 1.23, $p_a$ 14.7\,psia, length 11.518\,in,
their base model. \textbf{Our start line:} GENO's characteristic of Rao's field (GENO's lip Mach 1.7: at 1.5 one cell
of that fixed line is not certified).

\smallskip
\verdict{\HTHVERDICT}
\smallskip
\why{\HTHWHY}
\end{column}
\end{columns}
\end{frame}

% =====================================================================================================================
\caseframe{Bell, $\varepsilon=4$, length $4\,r_t$: Rao's thrust-optimised contour}{17_bell.png}{Rao 1958: maximum vacuum thrust at fixed length and exit radius (GENO type 2, CH$_4$/O$_2$
frozen gas, 54\,\% of the ideal length). Ours: throat arc to $\theta_B$, then a curve through 8 points, lip pinned;
reference = the best such curve of GENO's wall.}{num_bell_table}{\BELLVERDICT}{\BELLWHY}

% =====================================================================================================================
\caseframe{Bell, $\varepsilon=4$, at its ideal length: the ideal bell}{17_bellideal.png}{In vacuum at fixed exit area the thrust is at most $F_\mathrm{1D}=\dot mV_e+p_eA_e$, reached
only by a uniform parallel exit: GENO's ideal bell (type 0), length $7.34\,r_t$. Reference: its closest curve.}{num_bellideal_table}{\BELLIDEALVERDICT}{\BELLIDEALWHY}

% =====================================================================================================================
\caseframe{Annular nozzle (GENO's Migdal walls): carrier and production walker}{17_migdal.png}{Inlet and exit annulus fixed, vacuum: $C_F\le C_{F,\mathrm{1D}}=1.647385$, reached by
Migdal's perfect nozzle (GENO's walls). Both walls are designed: 12 + 12 points.}{num_migdal_table}{\MIGDALVERDICT}{\MIGDALWHY}

% =====================================================================================================================
\begin{frame}{When nobody knows the optimum: what the tool tells you}
\begin{columns}[T]
\begin{column}{0.5\textwidth}
\small
On a new design there is no Rao to compare with. The tool reports two numbers of its own, and the known cases
calibrate them:
\begin{enumerate}
\item \textbf{Stop reason and predicted remaining gain} $\tfrac12g^\top B^{-1}g/J$ with $B$ measured at the landing:
``CONVERGED'' means below $10^{-8}$.
\item \textbf{Spread between two starts}: run twice from distant walls; the two landings must agree in thrust and
wall.
\end{enumerate}
\plain{Rule of thumb.}{If both starts say CONVERGED and agree to $<1$\ppm, the landing is the optimum of the posed
problem up to grid and wall description (a few ppm). If they disagree, the problem has two tops or the wall
description is too coarse --- not a reason to trust either. The ideal-bell rows show why: one start says ``converged''
(0.006\ppm\ left to gain) on a top 29\ppm\ below the other's.}
\end{column}
\begin{column}{0.48\textwidth}
\centering\footnotesize
\resizebox{\linewidth}{!}{% generated by tests/t18_numbers.py -- do not edit
\begin{tabular}{@{}llrr@{}}\toprule
case, start & stop & predicted & vs closest curve\\
 & & remaining (ppm) & optimiser's grid (ppm)\\\midrule
Planar plug: 15\textdegree\ fan plug & converged & 0.003 & +1.70\\
Planar plug: straight line & converged & 2e-05 & +1.71\\
Methane plug: deformed Rao & converged & 0.001 & +1.95\\
Methane plug: straight line & converged & 0.002 & +1.95\\
Rao 1961: deformed Rao & converged & 0.0002 & +0.77\\
Rao 1961: Rao--line average & converged & 0.002 & +0.77\\
Axial exit: deformed Rao & converged & 0.0002 & +0.45\\
Axial exit: Rao--line average & converged & 0.008 & +0.44\\
Humphreys 1971: deformed Rao & converged & 6e-05 & +0.03\\
Humphreys 1971: Rao--line average & converged & 0.002 & +0.02\\
Bell TOC: truncated ideal & converged & 0.003 & +0.75\\
Bell TOC: squeezed TOC & converged & 0.004 & +0.75\\
Ideal bell: stretched TOC & converged & 0.006 & -28.8\\
Ideal bell: TOC + cylinder & converged & 0.001 & +0.77\\
Annular: no rule & converged & 0.0002 & +1.84\\
Annular: no-fold rule on & converged & 0.001 & +1.84\\
\bottomrule
\end{tabular}
}
\end{column}
\end{columns}
\end{frame}

% =====================================================================================================================
\begin{frame}{What is still open (1/2): precision and multiple tops}
\small
\begin{itemize}
\item \OPENONE
\item \OPENTWO
\end{itemize}
\end{frame}

\begin{frame}{What is still open (2/2)}
\small
\begin{itemize}
\item \OPENTHREE
\item \textbf{Physics, not code:} at low chamber pressure a plug designed for the mean state over-expands the flow;
compressions coalesce, no shock-free solution exists, and the 12-phase objective is undefined there.
\item \textbf{Not covered:} the walker's basis map and capped posings; the 12-phase family generation; base-pressure
closures; swirl; the inlet-family wrapper; the CFD referee. The constant-Mach (SDT) inflow is parked.
\item \textbf{Git:} nothing committed (owner's order). Changed tracked files: \texttt{a1\_plug\_spline\_opt.py} (driver
switch), \texttt{a1\_ideal\_march\_jax.py} (rescue default); new: \texttt{a1\_qntr.py}. Other sessions using the same tree
were told how to restore the old defaults. Keep/remove list to agree.
\end{itemize}
\end{frame}

% =====================================================================================================================
\appendix
\begin{frame}{Backup: the ideal bell's two tops}
\begin{columns}[T]
\begin{column}{0.5\textwidth}
\includegraphics[width=\linewidth]{../figs/23_bell_two_tops.png}

\scriptsize
\resizebox{\linewidth}{!}{\begin{tabular}{@{}lcc@{}}\toprule
 & second top & ideal bell\\\midrule
attachment angle & 8.16\textdegree & 12.53\textdegree\\
Hessian eigenvalues ($/J$) & $-18.7\dots-0.21$ & $-18.7\dots-0.34$\\
all 8 negative & \yes & \yes\\
true marches at $\pm2\times10^{-3}$, $\pm10^{-2}$ & 32/32 lower & 32/32 lower\\
gain still predicted & 0.015\ppm & 0.001\ppm\\\bottomrule
\end{tabular}}
\end{column}
\begin{column}{0.48\textwidth}
\scriptsize
\textbf{What is claimed.} With the 8-point wall on the $81\times240$ grid, the thrust the march computes has two
\emph{strict local maxima}, 30--45\ppm\ apart.

\smallskip
\textbf{The evidence} (t23, t24). At both results: the curvature (secant Hessian of the marched thrust) is negative in all
8 principal directions, and the true march, moved along each of them in both senses by $2\times10^{-3}$ and $10^{-2}$,
always gives less thrust (64 marches, largest change $-0.55$\ppm). Along the straight line between the two results the
thrust falls 14\ppm\ below the lower one before rising to the higher (left). So the driver's ``converged'' was right:
not a saddle, not a stop short.

\smallskip
\textbf{What is not settled.} Whether the second top exists in the continuous problem: it is not the throat-arc panel the
objective leaves out (the objective is smooth where the arc gains a station) nor the axial-Mach margin (0.04\,m/s at
both); with a 12-point wall the same start meets jumps of the march ($-230$ to $-345$\,N over wall changes of
$2\times10^{-5}$) and stops before either top, while the cylinder start finds the ideal bell again. Two distant starts
are the protection.
\end{column}
\end{columns}
\end{frame}

\begin{frame}{Backup: the references reproduce the papers}
\footnotesize
The optimum of the Rao cases is GENO's contour; here GENO and our solver of Rao's equations are tied to the papers
(the case slides show the papers' points together with the SQP results).

\footnotesize
Rao's exit-characteristic conditions solved from the paper alone (\texttt{ref/rao1961.py}):
$w\cos(\theta+\alpha)/\cos\alpha=C_2$, $R\rho w^2\sin^2\theta\tan\alpha=C_3$, corner conditions at $E$ ($p_a$) and $D$.

\smallskip
\resizebox{0.6\linewidth}{!}{\begin{tabular}{@{}lccc@{}}\toprule
$\varepsilon=3.81$ & Rao, printed & ours at $R_D$ & GENO at $R_D$\\\midrule
$L/R_E$ & 1.164 & 1.1640 & 1.1640\\
$\pi R_E^2/A^*$ & 3.81 & 3.8094 & 3.8095\\
$C_F$ (vacuum) & 1.5804 & 1.5804 & 1.5806\\
$\theta_D$ & $-19.73^\circ$ & $-19.69^\circ$ & $-19.69^\circ$\\\bottomrule
\end{tabular}}

\smallskip
Rao's tables = his equations ended at his printed $R_D$ (a base pressure of $0.0033\,P_c$, one march step early).
Humphreys 1971: $R_E=8.3311$ in (printed 8.33), $R_D=1.37507$ (1.37506), thrust 34\,254\,lbf (34\,253); GENO within
0.0003\,in of their 19 points; their $-12.27^\circ$ at $x=5.337$ is a misprint for $-22.27^\circ$.

\end{frame}

\begin{frame}{Backup: bells in the literature --- Allman \& Hoffman 1981}
\begin{columns}[T]
\begin{column}{0.5\textwidth}
\footnotesize
Allman \& Hoffman (AIAA J.\ 19(6), 1981) optimised bell walls \emph{directly} --- a 2nd-degree polynomial, three
search methods --- and compared with Rao's contours at 15 lengths: vacuum thrust 0.05--0.21\,\% (500--2100\ppm)
below Rao's, 0.22--0.66\,\% with ambient pressure; ``each method converges to the unique maximum thrust contour'', the
deficit being ``the restriction of a second-degree polynomial''.

\medskip
Here: an 8-point curve, the same kind of direct optimisation, lands within \BELLVSAH\ of the closest curve to Rao's: the
representation, not the method, sets the deficit --- their conclusion, one wall description later.
\end{column}
\begin{column}{0.48\textwidth}
\footnotesize
\begin{tabular}{@{}rrrr@{}}\toprule
length (cm) & Rao (N) & polynomial (N) & $\Delta F$\\\midrule
5.897 & 21362.4 & 21341.9 & $-0.10\,\%$\\
11.616 & 23315.1 & 23284.8 & $-0.13\,\%$\\
24.518 & 25217.5 & 25186.8 & $-0.12\,\%$\\
42.968 & 26402.0 & 26361.5 & $-0.15\,\%$\\
62.221 & 27050.1 & 26991.4 & $-0.21\,\%$\\\bottomrule
\end{tabular}

\smallskip
(their Table 1, vacuum rows; conditions in AFAPL-TR-77-1, not reproducible here)
\end{column}
\end{columns}
\end{frame}

\begin{frame}{Backup: findings on GENO (reference code, not modified)}
\footnotesize
\begin{itemize}
\item Rao's input $\theta_i$ is not the flow angle at the lip: GENO places $(M_E,\theta_E)$ by one midpoint step of
$\theta+\nu=$ const; inputs must emulate that step (done exactly here).
\item Rao cannot start from a user's characteristic: its inflow is an output; the model takes GENO's.
\item Default exit sampling $N_e=100$ makes polygon walls ($6\times10^{-3}$ rad per vertex): use $N_e\ge1000$.
\item The lip-fan scan covers only $60^\circ$; the ideal spike (control surface to the axis) diverges ($c^2v/y$ at
$y\to0$): the axial-exit case ends the plug at $R_D=0.4$.
\item TOP (type 4) at the ideal length reports $C_F$ above the ideal bell's ($+3.6\times10^{-5}$ in $I_{sp}$): a
parabola with an internal shock, credited isentropically --- not used as a reference here.
\item Type 3 (``conical'') with this input returns the ideal contour; the cone start was built analytically (and is
outside the shock-free model: a cell crosses the axis).
\end{itemize}
\end{frame}

\begin{frame}{Backup: reproduction}
\footnotesize
On s2 (venv \texttt{venv-a1}), from \texttt{RDE\_Design/validation/\_revalidation\_2026-10-07/}:
\begin{itemize}
\item walks: \texttt{run\_walk.sh <log> <cpu> <script> <args>} with \texttt{t08a\_planar\_walk.py},
\texttt{t08b\_rao\_walk.py}, \texttt{t08c\_migdal\_walk.py}, \texttt{t10\_rao\_walk.py}, \texttt{t15\_bell\_walk.py};
the production walker: \texttt{TWOP\_STAGE=walk TWOP\_START=F TWOP\_METRIC=start python a1\_twowall.py} (log
\texttt{logs/t13\_walker\_F\_qntr.log});
\item readings and figures: \texttt{t17\_slides.py <case>}, numbers \texttt{t18\_numbers.py}, then
\texttt{tectonic slides.tex} in \texttt{doc/};
\item GENO on s3: \texttt{run\_geno.sh <case>} (\texttt{geno/bell\_*}, \texttt{geno/rao61\_*}, \dots).
\end{itemize}
Driver: \texttt{PSPL\_DRIVER=qntr} (default) / \texttt{legacy}; tolerance \texttt{QNTR\_TOL} (default $10^{-8}$);
Newton rescue \texttt{A1\_NEWTON\_RESCUE=1} (default) / 0.
\end{frame}

\end{document}
